\documentclass[11pt]{article}
\usepackage[T1]{fontenc}
\usepackage[utf8]{inputenc}
\usepackage[margin=1in]{geometry}
\usepackage{hyperref}
\hypersetup{colorlinks=true, urlcolor=blue, linkcolor=blue}
\usepackage{parskip}
\pagestyle{empty}

\begin{document}

\noindent
Chun-Yeol You \\
Department of Physics and Chemistry, DGIST \\
(Daegu Gyeongbuk Institute of Science and Technology) \\
Daegu 42988, Republic of Korea \\
\href{mailto:cyyou@dgist.ac.kr}{cyyou@dgist.ac.kr}

\vspace{0.5em}
\noindent July 29, 2026

\vspace{1em}
\noindent To the Editor-in-Chief, \\
ACM Transactions on Mathematical Software

\vspace{1em}
\noindent\textbf{Re: Submission of the manuscript ``Tuning Internal Parameters
of Algebraic Multigrid: A Cross-Library Study of Achievable Speedup and
Cheap-Feature Prediction''}

\vspace{0.5em}
Dear Editor,

Please consider the enclosed manuscript for publication in ACM Transactions on
Mathematical Software as a research article.

Algebraic multigrid (AMG) is the workhorse preconditioner embedded in every
major sparse-solver framework, yet its performance depends strongly on internal
parameters --- coarsening scheme, smoother, and strength-of-connection
threshold --- that most users leave at library defaults. Prior
algorithm-selection work almost exclusively chooses \emph{among} solvers or
preconditioners, treating an entire AMG method as a single label. Our manuscript
instead opens that label and studies the interior parameter space empirically:
we sweep 88 configurations of AMGCL over 150 (later 300) SuiteSparse matrices
--- 13{,}200 timed solves --- and ask how much speed is lost by not tuning, and
whether a good configuration can be predicted from features that are cheap
relative to the solve itself.

The main findings are:

\begin{itemize}
\item \textbf{Tuning matters and requires per-matrix selection.} The oracle
speedup over the default has median $2.4\times$ (upper quartile $6.4\times$,
maximum $292\times$); the best configuration varies across 14 distinct
coarsening/smoother labels, and on 31 matrices the default fails outright while
some configuration succeeds.
\item \textbf{A predictor works from near-free features, evaluated honestly.}
Under a strict leave-one-group-out protocol --- no sibling matrices leaking
across the split, a check the literature commonly omits --- the predictor picks
a configuration that solves 83--89\% of held-out matrices and captures a median
43--49\% of the achievable speedup. Near-free structural features suffice;
spectral features costing $100$--$1000\times$ more add nothing, a feature-cost
accounting that prior work seldom performs.
\item \textbf{A self-critical cross-library validation.} Repeating the study on
hypre BoomerAMG shows that tunability replicates qualitatively but its magnitude
is smaller (median $1.6\times$) and does not transfer between libraries (rank
correlation $0.24$); it is largely governed by default quality --- the dramatic
AMGCL speedups occur precisely where hypre's well-engineered default already
wins. Intrinsic solvability, by contrast, is library-independent (hard-matrix
Jaccard $0.76$), a split that a predictor-transfer experiment confirms in both
directions.
\end{itemize}

We believe the work fits TOMS well on both scope and values. It is an empirical
study of two widely-used pieces of mathematical software (AMGCL and hypre), it
foregrounds honest methodology (feature-cost accounting and group-wise
evaluation), and it is fully reproducible. We release the complete artifact ---
the 18{,}600-solve dataset, the feature extractor, the predictor, and all
figure-generation scripts --- at
\url{https://github.com/mirryou-maker/amg-tuning-crosslib} (currently private;
it will be made public upon acceptance, and we are glad to provide reviewer
access immediately). A single command, \texttt{python reproduce.py}, regenerates
every figure and headline number from the released data, and a companion tool,
\texttt{tools/recommend.py}, applies the predictor to a user's own matrix. We
would welcome consideration under the Replicated Computational Results process.

We confirm that this manuscript is original, has not been published previously,
and is not under consideration elsewhere. In accordance with the ACM Policy on
Authorship, the manuscript includes a Generative AI disclosure describing the
use of an AI coding assistant (Claude Code, Anthropic) for implementation,
experiment orchestration, and drafting; the author verified all experimental
design, numerical results, and claims against the released code and data, and
takes full responsibility for the content. The work was supported by the
National Research Foundation of Korea (Nos.\ RS-2025-25463492 and
RS-2026-25472340) and the DGIST R\&D Program (26-SR-01).

\textbf{Suggested reviewers.} We suggest the following experts in algebraic
multigrid, solver selection, and numerical software (contact details available
on request); we have no conflict of interest with them:

\begin{itemize}
\item \textbf{Luke N.\ Olson}, University of Illinois Urbana-Champaign ---
algebraic multigrid and the PyAMG software.
\item \textbf{Jed Brown}, University of Colorado Boulder --- multigrid, PETSc,
and numerical-software performance.
\item \textbf{Kanika Sood}, California State University, Fullerton --- machine
learning for iterative-solver selection (the closest prior line of work).
\item \textbf{Scott P.\ MacLachlan}, Memorial University of Newfoundland ---
multigrid methods and their analysis.
\item \textbf{Sivasankaran Rajamanickam}, Sandia National Laboratories ---
sparse linear algebra, performance, and learning-based methods for solvers.
\end{itemize}

Thank you for considering our submission.

\vspace{1em}
\noindent Sincerely,

\vspace{1.5em}
\noindent Chun-Yeol You \\
Department of Physics and Chemistry, DGIST \\
\href{mailto:cyyou@dgist.ac.kr}{cyyou@dgist.ac.kr}

\end{document}
